\chapter{Problem analysis}
Regarding the name of the project, we decided to hold a convention laid down by the previous project made by the PeachPie developers, where the objective was to compile Wordpress framework into .NET. The resulting project was named Peachpied Wordpress, therefore, this project was entitled Peachpied Symfony and all custom tools provided use the Peachpied.Symfony prefix in their names.

\paragraph{Symfony project structure} As was mentioned before, each Symfony project is made from libraries, generally called Symfony Components. All these libraries reside in the \texttt{vendor} directory of the project. The rest of the project contains mainly settings such as, routing definition, environmental variables or database credentials configuration files, and user defined content, such as controller classes responsible for handling requests, images, definition of styling and front-end scripts. Aside from this, every Symfony project also contains a set of autoload files. Those are responsible for the proper load-up of all of the remaining libraries included in the vendor folder and are stored under the \texttt{vendor} directory.

\paragraph{Current problems} Ability to compile Symfony projects into .NET would yield undeniable contributions to both the communities. In terms of plain feasibility, this is what Peachpie compiler already does, nevertheless, aside from its dependence on using PHP and composer to generate certain files, there are also other problems that recede this approach from being usable in practise. These will be discussed right ahead.

When using plain Peachpie to compile Symfony project, the major problems lay in efficiency. Each time Symfony project is compiled, the whole content of the \texttt{vendor} folder is recompiled. For a bare bones project this means over a 15MB of \texttt{.php} code sources distributed over 3000+ files. This equals tenths of minutes of compilation each time. And we are talking only about the bare bones project with approximately 30 Symfony Components. The actual projects may use hundreds of them making this approach impossible in practise. We must not forget the disk space consumed by the resulting \texttt{.dll} which would be appropriate to the size of the project. Finally, due to differences in type systems and execution models of PHP and .NET, compilation errors calling either for some files' exclusion from the compilation or their direct modifications are to be expected as well.

Aside from using Peachpie in the previously mentioned manner to embed an existing Symfony project into our new .NET Core application, there is also a way of using Peachpie to include only a certain functionality from Symfony. This can be done by embedding only a set of Symfony Components in our application. This approach seems more promising as there would definitely be less files to be compiled and thus less compilation errors that needs to be tend to. Nevertheless, these files would require to be compiled over and over as well. Furthermore, there would certainly be an urge to write an API for using the desired functionality from .NET which would require a non-trivial knowledge of both Peachpie and the Symfony Component.


\paragraph{Aim of the project} When measured by the complexity of the task in terms of time needed for its completion or spatial and performance overheads, Peachpie compiler would definitely need a help of more specialized tool to deliver the solution in practically usable form. The aim of this project is to address this need and reduce both compilation time of mentioned ways of embedding Symfony within .NET Core application while increasing its actual usability by providing simple API for its usage. The spacial requirements were also considered but did play only secondary role during proposed solution's design. 



\section{Proposed Solution}
Taking its inner structure into account, one can perceive the Symfony project as two completely distinct parts. First, the \texttt{vendor} directory containing all of the libraries that the project uses, and second, basically the rest of the project folder. Moreover, for the set of Symfony components and fixed versions, the vendor folder will always look the same. Ultimately, each of these Symfony Components can be recognized as an independent unit as well, only holding references to another libraries that it requires for its proper functioning via its own \texttt{composer.json} file.

The whole solution of the previously mentioned problem is therefore based on these conclusions:

\begin{enumerate}[1:]
\item User does not manually modify content of the vendor directory.
\item The \texttt{vendor} directory for a set of Symfony Components with fixed versions will be always identical (excluding the autoloaders part). 
\item Symfony Components are self sustained with functionalities depending only on libraries listed in included \texttt{composer.json} file.
\end{enumerate}


\paragraph{Component compilation} Taking this inner composition into account, we decided that pre-compiling Symfony components and proving each one on demand would bring the most desirable outcome. This step will save considerable amount of time by deferring to the pre-compilation while bringing the highest flexibility regarding future extensibility. NuGet package manager was selected to handle this as a most widespread package manager under .NET environment. This way each composer dependency will be simulated by a package reference and NuGet will handle the rest while also saving physical space thanks to its caching system.

One drawback of introduced approach is the necessity of significant number of packages to be pre-compiled in order to comply the requirements even of the bare bones Symfony 4.2.3 project, created with the \texttt{composer create-project Symfony/Skeleton [projectName] [<version>]} command. Such project is called Symfony Skeleton and consists of 24 components. There is no feasible solution to this other than providing an at least partially automatized way of producing packages for Symfony Components.

\paragraph{Project compilation} With this approach, we will be able to supply Symfony project with pre-compiled libraries in a form of NuGet packages of required versions while saving both the compilation time and the size of resulting \texttt{.dll}. This also does resolve some of the previous dependence on PHP and composer, however, it brings several new questions that has to be addressed in order to move forward. How to handle composer's autoloading? What about Symfony's caching directory containing \texttt{.php} files produced on application start-up? And how will be composer's dependencies and the \texttt{composer.lock} file mapped to package references?

Autoloading scripts in Symfony are produced by composer's \texttt{dump-autoload} or \texttt{install} command and updated each time new Symfony Component is installed. It is true that their content is determined by a set of used components and their versions, however, we can see that providing them in a pre-compiled state as NuGet packages as well would prove very challenging. Choosing this approach would require to produce and store NuGet packages for autoloading scripts of all possible combinations of packages and their compatible versions. One, and probably the best, alternative to this approach is using composer itself to produce these scripts during an early phase of project's compilation and then simply compile them with the project. Time overhead of this method is negligible, compared to time needed for autoloading scripts pre-compiling.

Symfony projects use caching mechanism in order to speed up future request handling. These application cache are by default stored in the \texttt{var} directory and produced after first application startup. In order for the application to work correctly under .NET, these files need to be compiled as well. For Peachpie however, this means the application has to be started first either from PHP or from .NET after previous compilation, and then recompiled again with generated cache files. An easy and efficient solution to this lays once again in producing these files manually during an early phase of of project's compilation.

On he contrary of previous conclusions, solution to the final question of how to convert composer's dependencies to package references might lay already within Peachpie. One of the targets provided by the Peachpie SDK parses \texttt{composer.json} file and converts every dependency into package reference. Unfortunately, it does so only for packages not marked as development dependencies. In addition to, it does not check whether does a NuGet package for listed dependency exist in any of registered repositories. This is problematic because in .NET every package used somewhere in the code needs to be referred and thus cannot be considered a development dependency. Moreover, some packages non-development components provide only additional CLI (Command Line Interface) functionalities for composer and PHP, and therefore do not need to be referred nor compiled. Correct way to handle this issue would be to cooperate with Peachpie developers and adjust current tool to befit out needs. We decided to avoid this way in sake of saving the time such task would consume and postpone it into possible future works. Writing a custom tool handling this seems as a quick alternative, nevertheless it would require delivering it before NuGet's \texttt{restore} target is executed which means loading it directly from a \texttt{.dll}. From the perspective of UX (User Experience), this approach brings unpleasant complications, and after comparing it with complexity of setting up package references manually, we decided to settle for the manual solution for now. 


\paragraph{} Its is clear now, that regarding the number of Symfony Components and continuing development of Peachpie, delivering our tool in a production grade quality would require tremendous amount of time or a greater team of developers. Because that exceeds both our capabilities and the scope of the thesis we focused primarily on outlining how the tools should work and displaying some possibilities to provide further interoperability. The whole project is therefore to be regarded as a proof of concept.
